El proyecto se restringe al flujo de creación individual de materiales dentro de SAP, que es el mecanismo mediante el cual un gestor solicita, valida y registra un material nuevo en el catálogo de la unidad de negocio. Dentro de este flujo, GAMMA interviene en tres momentos: normaliza la descripción propuesta por el gestor, verifica su similitud contra el catálogo cargado en la plataforma para alertar sobre posibles duplicados, y sugiere la categoría más probable con base en el histórico de materiales ya clasificados. GAMMA es una herramienta de asistencia y validación: no escribe directamente en SAP ni sustituye el criterio del gestor, quien conserva el control de la decisión final. Las solicitudes ya validadas se entregan mediante exportación a un archivo de hoja de cálculo, que es el insumo con el que el material se registra posteriormente en el sistema.

La plataforma no mantiene integración en línea con el ERP, ni de lectura ni de escritura. Esta separación es deliberada y constituye el mecanismo de supervisión humana sobre el que se diseñó la solución: el gestor exporta periódicamente el maestro de materiales desde SAP hacia la plataforma, y registra en SAP las solicitudes que la plataforma ya validó, de modo que ninguna operación sobre el sistema transaccional ocurre sin la intervención de una persona. Su contrapartida es que la detección de duplicados alcanza a los materiales existentes hasta la fecha de la última carga, por lo que un registro creado en SAP con posterioridad no es considerado en la comparación hasta que el catálogo se actualice.

El modelo de clasificación se entrena y valida sobre el catálogo de una única unidad de negocio, que es el ámbito de esta primera implementación. Su extensión al resto de las unidades de la corporación queda fuera del alcance del proyecto y se plantea como la etapa siguiente, una vez confirmado su desempeño en este entorno; dentro del catálogo empleado, las categorías con muy escasa representación histórica quedan fuera del espacio de predicción del modelo. El modelo se despliega, asimismo, como un artefacto entrenado de forma previa: la incorporación de un proceso automatizado de reentrenamiento que lo realimente con las correcciones registradas por los gestores durante el uso queda fuera del alcance de esta fase y se identifica como la evolución técnica prevista para la herramienta.

Queda fuera del alcance del proyecto el flujo de creación masiva de materiales, que opera mediante un mecanismo distinto al de la creación individual y no fue intervenido. Tampoco forma parte del alcance la corrección automática de los materiales duplicados o mal categorizados que ya existían en el catálogo antes de la implementación: el proyecto cuantifica ese universo mediante la auditoría retroactiva descrita en los objetivos, pero su remediación corresponde a un esfuerzo de limpieza de datos separado, fuera del alcance de este trabajo.

El desarrollo y la validación del proyecto se realizaron con datos y usuarios de una unidad de negocio específica dentro de la organización, con un período de uso activo por parte del equipo de gestores comprendido entre el 18 de junio y el 21 de julio de 2026. Los resultados de tiempo de gestión, calidad del catálogo y satisfacción de usuario que se presentan en este documento corresponden a ese alcance y a ese período.

El proyecto se desarrolla y se limita al alcance de una única unidad de negocio dentro de una organización con presencia en múltiples países de Centroamérica y el Caribe; no se extiende a otras unidades, áreas o sistemas que la organización pudiera operar de forma independiente.
